\documentclass{article}
\usepackage{iclr2026_conference,times}

% Math commands for spectral_partition_router_v2_2

\newcommand{\R}{\mathbb{R}}
\newcommand{\E}{\mathbb{E}}
\newcommand{\Prob}{\mathbb{P}}
\newcommand{\calD}{\mathcal{D}}
\newcommand{\calE}{\mathcal{E}}
\newcommand{\calN}{\mathcal{N}}
\newcommand{\Sphere}{\mathbb{S}}
\newcommand{\Cov}{\mathrm{Cov}}
\newcommand{\tr}{\mathrm{tr}}
\newcommand{\norm}[1]{\left\lVert #1 \right\rVert}
\newcommand{\newterm}[1]{\textbf{#1}}


\usepackage{amsmath,amssymb,mathtools}
\usepackage{booktabs}
\usepackage{enumitem}
\usepackage{hyperref}
\usepackage{url}

\title{Static Feature-Level Specialization under Spectrally Isotropic Representations}

\author{Anonymous Authors}

\begin{document}
\maketitle

\begin{abstract}
We study a controlled static routing problem for mixture-of-experts models.  The setting is deliberately simple: route-relevant features have uniform feature-family frequencies and globally isotropic spectrum.  The goal is not merely to equalize expert counts, but to obtain feature-level specialization: each expert should receive a balanced and coherent subset of feature families, and this specialization should help a controlled static learning objective when the target depends on feature identity.  We first show that spectral isotropy alone is insufficient: if feature labels are independent of an isotropic representation, no deterministic feature-only router can recover specialization.  We therefore formulate a middle regime in which globally isotropic features contain locally recoverable feature families.  Under this regime we propose Balanced Feature-Specializing Routing (BFSR): spectrally normalize the calibration features, recover feature families by unsupervised clustering, assign recovered families to experts under a balanced constraint, and route new tokens by recovered family identity.  The mass mechanism is family-frequency balance rather than projection quantiles.  Quantile routing remains a baseline that can equalize mass without necessarily preserving feature identity.  We give the static problem definition, theorem sketches, and an experiment protocol with positive and negative anchors.  Training-time preservation and dynamic router updates are outside the scope of this paper.
\end{abstract}

\section{Introduction}

A common load-balancing objective for mixture-of-experts (MoE) routing asks each expert to receive a similar number of tokens.  This is not the same as feature specialization.  A router can be perfectly balanced by assigning tokens in the current batch while ignoring stable feature identity.  Conversely, a router may discover meaningful features but route them unevenly if the feature distribution is imbalanced.  To make the clean case precise, we separate the static dispatch problem from later training dynamics.

This paper studies the following question:

\begin{quote}
If route-relevant features are globally spectrally isotropic and feature families occur with uniform frequency, can a static non-cheating router assign feature families to experts in a balanced and specialized way?
\end{quote}

The answer cannot follow from spectral isotropy alone.  A standard Gaussian representation with random independent labels is spectrally isotropic, but no feature-only router can infer those labels.  Thus the causal condition for specialization must include more than global isotropy.  We argue that the right controlled regime is:

\begin{center}
\fbox{\parbox{0.88\linewidth}{
\centering
\textbf{global spectral isotropy} + \textbf{uniform feature-family frequencies} + \textbf{locally recoverable feature families}.
}}
\end{center}

In this regime, spectral isotropy provides a clean metric; uniform family frequency makes balanced expert usage possible by balanced family assignment; and local recoverability makes feature-level specialization identifiable.

The resulting method is \newterm{Balanced Feature-Specializing Routing} (BFSR).  BFSR first normalizes the representation spectrum, then recovers feature-family prototypes, then assigns those recovered families to experts with a balanced constraint.  The router for a new token predicts its feature family and sends it to the assigned expert.  This mechanism is more causally tied to the clean-case assumptions than quantile-first routing: quantiles can equalize mass for many distributions, including ones with no recoverable feature structure, whereas BFSR uses feature-family frequency and recoverability directly.

\paragraph{Contributions.}
We make four contributions.
\begin{enumerate}[leftmargin=*]
    \item We define static good dispatch as balanced expert mass plus feature-level specialization and controlled static utility.
    \item We prove an impossibility anchor: spectral isotropy alone cannot imply feature specialization.
    \item We propose BFSR, a non-cheating static router for globally isotropic but locally structured feature distributions.
    \item We give a falsifiable experiment plan that distinguishes mass balance, feature specialization, spectral normalization, and static utility.
\end{enumerate}

\paragraph{Scope.}
This paper does not study representation drift, route preservation, slow updates, hysteresis, or trainable router dynamics.  Those are later-stage problems.  The present goal is to close the static clean case.

\section{Controlled Static Setting}

Let \(X\in\R^d\) be a route-relevant representation, \(F\in\{1,\ldots,K\}\) a latent feature family, and \(E\) the number of experts.  A static router is a deterministic function
\[
    g:\R^d\to\{1,\ldots,E\}.
\]
We write \(G=g(X)\) for the selected expert.

\subsection{Feature-family distribution}

The clean controlled distribution is
\begin{equation}
    F\sim\Unif(\{1,\ldots,K\}),
    \qquad
    X\mid F=k = \mu_k + \epsilon,
    \label{eq:mixture}
\end{equation}
where the aggregate representation is globally isotropic after centering or whitening:
\begin{equation}
    \Cov(\widetilde X) \approx I_d.
    \label{eq:isotropy}
\end{equation}
This model allows local feature structure and global spectral uniformity to coexist.  For example, the centers \(\{\mu_k\}\) may form a balanced spherical code, while each family is locally concentrated around its center.

\subsection{Recoverability}

A feature family is useful for routing only if it is recoverable from the representation.  We assume that there exists an unsupervised feature-family estimator \(\widehat h\) such that, up to a permutation of family names,
\begin{equation}
    \Prob(\widehat h(X)\neq F)\leq \eta.
    \label{eq:recovery}
\end{equation}
This is a genuine additional assumption beyond spectral isotropy.  Without it, feature specialization is not identifiable.

\subsection{Non-cheating rule}

Router construction may use unlabeled calibration features, feature mean and covariance, unsupervised clustering, estimated cluster masses, and optional model-internal expert weights.  It may not use latent feature labels, task labels, next-token answers, evaluation labels, or current-batch expert counts.  Synthetic labels are used only to generate data and evaluate specialization.

\section{What Is Static Good Dispatch?}

A static router is good in this paper if it satisfies three objectives.

\paragraph{Balanced expert mass.}
Each expert should receive equal probability mass:
\begin{equation}
    \Prob(G=e)\approx \frac1E.
    \label{eq:mass}
\end{equation}
This is necessary but not sufficient.

\paragraph{Feature-level specialization.}
The router should preserve feature-family identity.  One form is small conditional entropy:
\begin{equation}
    \Entropy(G\mid F) \quad \text{small}.
\end{equation}
Equivalently, NMI, purity, or mutual information between \(G\) and \(F\) should be high when the feature families are recoverable.

\paragraph{Static utility.}
Under a controlled feature-dependent target, a specialized fixed router should improve expert learning relative to balanced but incoherent routing.  A typical regression target is
\begin{equation}
    Y = a_F^\top X + \xi.
    \label{eq:target}
\end{equation}
If \(Y\) is independent of feature identity in the relevant sense, feature specialization should not be expected to help.

\section{Why Spectral Isotropy Alone Is Insufficient}

Global isotropy does not define feature identity.  The following null anchor makes this precise.

\paragraph{Proposition 1: isotropy-alone non-identifiability.}
Let
\[
    X\sim\Normal(0,I_d),\qquad
    F\sim\Unif(\{1,\ldots,K\}),\qquad
    F\perp X.
\]
Then for any deterministic router \(g(X)\),
\begin{equation}
    \MI(g(X);F)=0.
\end{equation}

\paragraph{Proof.}
Since \(g(X)\) is measurable with respect to \(X\), and \(F\perp X\), we have \(F\perp g(X)\).  Hence the mutual information is zero.
\hfill\(\square\)

This proposition prevents a common overclaim: spectrally isotropic features do not by themselves imply good routing.  They provide a clean global metric only when feature families are also locally recoverable.

\section{Balanced Feature-Specializing Routing}

BFSR is a static router built from four steps.

\subsection{Step 1: spectral normalization}

Given calibration features \(\{x_i\}_{i=1}^n\), estimate
\[
    \widehat\mu = \frac1n\sum_i x_i,
    \qquad
    \widehat\Sigma = \frac1n\sum_i (x_i-\widehat\mu)(x_i-\widehat\mu)^\top.
\]
Then normalize
\begin{equation}
    \widetilde x_i = (\widehat\Sigma + \lambda I)^{-1/2}(x_i-\widehat\mu).
\end{equation}
This is a metric repair step, not a mass-balancing theorem.

\subsection{Step 2: recover feature families}

Fit an unsupervised family estimator
\[
    \widehat h:\R^d\to\{1,\ldots,K\}
\]
on normalized calibration features.  In the first controlled suite, \(K\) is known from data construction and \(\widehat h\) can be implemented by k-means or spherical k-means.

Let \(\widehat\mu_c\) be the recovered family prototype and \(\widehat p_c\) its estimated mass.

\subsection{Step 3: balanced family-to-expert assignment}

Find an assignment
\[
    a:\{1,\ldots,K\}\to\{1,\ldots,E\}
\]
subject to
\begin{equation}
    \sum_{c:a(c)=e}\widehat p_c \approx \frac1E.
    \label{eq:assignment_balance}
\end{equation}
When family frequencies are uniform and \(K\) is divisible by \(E\), the simplest assignment gives exactly \(K/E\) families to each expert.

The assignment can be unsupervised or expert-aligned.  An unsupervised locality objective groups nearby recovered families:
\begin{equation}
    \min_a \sum_{e=1}^E \sum_{c:a(c)=e} \widehat p_c
    \norm{\widehat\mu_c-\bar\mu_e}^2
\end{equation}
under the balance constraint.  If expert prototypes \(u_e\) exist, an expert-aligned objective is
\begin{equation}
    \max_a \sum_{c=1}^K \widehat p_c\, s(c,a(c)),
\end{equation}
where \(s(c,e)\) is cosine or projective similarity.  Projective similarity avoids sign ambiguity:
\[
    P_c = \frac{\widehat\mu_c\widehat\mu_c^\top}{\norm{\widehat\mu_c}^2},\qquad
    P_e = u_eu_e^\top,\qquad
    s(c,e)=\tr(P_cP_e).
\]

\subsection{Step 4: route new tokens}

For a new token feature \(x\), normalize it, predict its recovered family \(\widehat h(\widetilde x)\), and route by
\begin{equation}
    g(x)=a(\widehat h(\widetilde x)).
\end{equation}

\subsection{Variants and controls}

\begin{center}
\begin{tabular}{lll}
\toprule
Method & Uses labels? & Role \\
\midrule
BFSR-OracleFeature & yes & invalid upper bound \\
BFSR-Cluster & no & primary non-cheating method \\
BFSR-ExpertAligned & no task labels & tests expert-side geometry \\
BFSR-RandomGroup & no & balanced feature grouping control \\
Quantile-SPR & no & mass-only baseline \\
BalancedAssignment & current batch counts & invalid distributional control \\
\bottomrule
\end{tabular}
\end{center}

\section{Theory}

\paragraph{Theorem 2: balanced family assignment gives balanced expert mass.}
Assume \(F\sim\Unif(\{1,\ldots,K\})\), \(K\) is divisible by \(E\), and the ideal assignment \(a\) maps exactly \(K/E\) families to each expert.  If \(g(X)=a(F)\), then
\[
    \Prob(g(X)=e)=1/E
\]
for every expert \(e\).

\paragraph{Proof.}
\[
\Prob(g(X)=e)=\sum_{k:a(k)=e}\Prob(F=k)=\frac{K/E}{K}=\frac1E.
\]
\hfill\(\square\)

\paragraph{Theorem 3: robust mass under recovery error.}
Let \(a(F)\) be a balanced ideal family assignment.  Suppose the recovered family estimator satisfies \(\Prob(\widehat h(X)\neq F)\leq\eta\).  Define \(g(X)=a(\widehat h(X))\).  Then for every expert \(e\),
\begin{equation}
    \abs{\Prob(g(X)=e)-1/E}\leq \eta.
\end{equation}

\paragraph{Proof sketch.}
The ideal expert variable \(G^\star=a(F)\) is balanced by Theorem 2.  The implemented expert variable \(G=a(\widehat h(X))\) differs from \(G^\star\) only on events where \(\widehat h(X)\neq F\).  Hence the total variation distance between the two expert distributions is at most \(\eta\), which implies the coordinate-wise bound.
\hfill\(\square\)

\paragraph{Theorem 4: specialization under recovery error.}
Under the same assumptions,
\begin{equation}
    \Prob(G\neq a(F))\leq \eta.
\end{equation}
Thus feature-level specialization degrades with feature recovery error.  Conditional entropy or mutual-information versions follow from standard error-to-information inequalities.

\paragraph{Proposition 5: utility requires feature-dependent structure.}
If the target is feature-independent and the expert class has no family-conditioned approximation advantage, feature-specializing dispatch should not improve Bayes risk.  If the target is family-dependent, e.g.
\[
    Y=a_F^\top X+\xi,
\]
and experts are capacity-limited, routing coherent feature families to experts can reduce approximation error relative to incoherent partitions that mix incompatible family-specific rules.

This proposition defines an experimental anchor rather than a universal theorem about all expert classes.

\section{Experiments}

This section is intentionally written as a protocol.  Numeric results and analysis should be filled after the experiment suite is run.

\subsection{Datasets}

\paragraph{D1: isotropic null.}
\[
    X\sim\Normal(0,I_d),\qquad F\sim\Unif(\{1,\ldots,K\}),\qquad F\perp X.
\]
Expected: valid non-cheating methods should have near-zero NMI with \(F\).

\paragraph{D2: isotropic feature-family mixture.}
\[
    F\sim\Unif(\{1,\ldots,K\}),\qquad X\mid F=k=\mu_k+\epsilon.
\]
Centers are balanced so the aggregate spectrum is isotropic after normalization.  Separation and noise are swept.

\paragraph{D3: anisotropically distorted mixture.}
Start from D2 and apply \(X_{\rm raw}=AX\) with controlled condition number.  Compare raw, centered, and whitened routing.

\paragraph{D4: frequency-imbalanced mixture.}
Use non-uniform \(\Prob(F=k)\).  This tests whether uniform frequency is the mass mechanism.

\paragraph{D5: static utility targets.}
Use both feature-dependent targets \(Y=a_F^\top X+\xi\) and feature-independent targets.

\subsection{Experiment list}

\begin{center}
\begin{tabular}{lll}
\toprule
ID & Question & Result status \\
\midrule
E1 & Clean mass and specialization sanity & TODO \\
E2 & Isotropy-alone null anchor & TODO \\
E3 & Feature recoverability sweep & TODO \\
E4 & Spectral normalization ablation & TODO \\
E5 & Frequency uniformity ablation & TODO \\
E6 & Expert alignment ladder & TODO \\
E7 & Static utility audit & TODO \\
\bottomrule
\end{tabular}
\end{center}

\subsection{Metrics}

Mass metrics include LoadCV, max mass error, and expert usage entropy.  Feature-specialization metrics include NMI, purity, \(\Entropy(G\mid F)\), \(\MI(G;F)\), feature split rate, and within-expert feature diversity.  Feature-recovery metrics include clustering accuracy up to permutation, ARI, and cluster purity.  Static utility metrics include validation loss and per-feature validation loss.

\subsection{Expected interpretation}

The experiments should support or falsify the causal map:
\begin{center}
\begin{tabular}{ll}
\toprule
Condition & Expected effect \\
\midrule
spectral isotropy alone & not enough for specialization \\
uniform feature frequencies & mass balance under balanced family assignment \\
local feature recoverability & high route-feature consistency \\
spectral normalization & metric repair under anisotropic distortion \\
feature-dependent target & static utility from specialization \\
\bottomrule
\end{tabular}
\end{center}

\paragraph{Results placeholder.}
All quantitative tables are intentionally left empty until experiments are run:
\begin{itemize}[leftmargin=*]
    \item Table 1: distribution diagnostics.
    \item Table 2: static dispatch metrics.
    \item Table 3: feature recovery metrics.
    \item Table 4: static utility metrics.
\end{itemize}

\section{Discussion}

The central distinction is between mass balance and feature specialization.  Quantile routing is a strong mass-control baseline, but the current clean case contains a more direct mass mechanism: if feature families are equally frequent, assigning equal numbers or equal total masses of families to each expert balances expert usage.  This assignment also preserves feature identity if the families are recoverable.

Spectral isotropy should also be interpreted carefully.  It is not the source of semantic or feature identity.  It removes global anisotropic confounds and defines a clean geometry in which local family recovery can be tested.  The paper's positive claim must therefore condition on local recoverability.

\section{Limitations}

This paper is limited to synthetic or controlled feature-family distributions.  It does not claim that real language-model hidden states naturally satisfy these assumptions.  It does not solve non-uniform spectra, feature-frequency imbalance beyond the specific ablation, dynamic route preservation, or learned-router training.  It also assumes that the number of feature families \(K\) is known in the first clean suite; estimating \(K\) is a separate problem.

\section{Conclusion}

We formulated a static clean-case problem for feature-level MoE specialization under globally spectrally isotropic and frequency-uniform feature distributions.  Spectral isotropy alone is insufficient, so the solvable regime must include locally recoverable feature families.  BFSR uses the resulting causal structure directly: normalize the metric, recover feature families, assign families to experts under a balanced constraint, and route by recovered family identity.  This closes the static problem only if experiments show predictable success under recoverability and predictable failure when recoverability or feature-dependent utility is removed.

\bibliography{references}
\bibliographystyle{iclr2026_conference}

\end{document}
